\section[鄭風]{鄭風}

\subsection[緇衣]{緇衣}

緇衣之宜兮，敝，予又改爲兮。適子之館兮，還，予授子之粲兮。

緇衣之好兮，敝，予又改造兮。適子之館兮，還，予授子之粲兮。

緇衣之蓆兮，敝，予又改作兮。適子之館兮，還，予授子之粲兮。

\subsection[將仲子]{將仲子}

將仲子兮，無踰我里，無折我樹杞。豈敢愛之？畏我父母。仲可懷也，父母之言，亦可畏也。

將仲子兮，無踰我牆，無折我樹桑。豈敢愛之？畏我諸兄。仲可懷也，諸兄之言，亦可畏也。

將仲子兮，無踰我園，無折我樹檀。豈敢愛之？畏人之多言。仲可懷也，人之多言，亦可畏也。

\subsection[叔于田]{叔于田}

叔于田，巷無居人。豈無居人，不如叔也。洵美且仁。

叔于狩，巷無飲酒。豈無飲酒，不如叔也。洵美且好。

叔適野，巷無服馬。豈無服馬，不如叔也。洵美且武。

\subsection[大叔于田]{大叔于田}

叔于田，乘乘馬。執轡如組，兩驂如舞。叔在藪，火烈具舉，襢裼暴虎，獻于公所。將叔無狃，戒其傷女。

叔于田，乘乘黄。兩服上襄，兩驂鴈行。叔在藪，火烈具揚。叔善射忌，又良御忌，抑磬控忌，抑縱送忌。

叔于田，乘乘鴇。兩服齊首，兩驂如手。叔在藪，火烈具阜。叔馬慢忌，叔發罕忌，抑釋掤忌，抑鬯弓忌。

\subsection[清人]{清人}

清人在彭，駟介旁旁。二矛重英，河上乎翱翔。

清人在消，駟介麃麃。二矛重喬，河上乎逍遥。

清人在軸，駟介陶陶。左旋右抽，中軍作好。
